% preamble.tex — pacotes e estilos do Guia da Secretaria
\usepackage[T1]{fontenc}
\usepackage[utf8]{inputenc}
\usepackage[brazil]{babel}
\usepackage{lmodern}
\usepackage{microtype}
\usepackage{geometry}
\usepackage{graphicx}
\usepackage{booktabs}
\usepackage{longtable}
\usepackage{array}
\usepackage{tabularx}
\usepackage{multirow}
\usepackage{xcolor}
\usepackage{tikz}
\usetikzlibrary{arrows.meta, positioning, shapes.geometric, fit, backgrounds}
\usepackage{hyperref}
\usepackage{fancyhdr}
\usepackage{titlesec}
\usepackage{enumitem}
\usepackage{tcolorbox}
\tcbuselibrary{skins, breakable}

\usepackage{ragged2e}

\geometry{a4paper, left=2.5cm, right=2.5cm, top=2.8cm, bottom=2.5cm}

% Colunas reutilizáveis para tabelas
\newcolumntype{Y}{>{\RaggedRight\arraybackslash}X}
\newcolumntype{L}[1]{>{\RaggedRight\arraybackslash}p{#1}}
\newcolumntype{C}[1]{>{\Centering\arraybackslash}p{#1}}

\newcommand{\guiaTableStretch}{\renewcommand{\arraystretch}{1.35}}

\definecolor{seplanblue}{RGB}{0, 82, 136}
\definecolor{seplanlight}{RGB}{230, 242, 250}
\definecolor{repgreen}{RGB}{22, 101, 52}
\definecolor{revblue}{RGB}{30, 64, 175}
\definecolor{warnamber}{RGB}{180, 83, 9}
\definecolor{tipcyan}{RGB}{14, 116, 144}

\hypersetup{
  colorlinks=true,
  linkcolor=seplanblue,
  urlcolor=seplanblue,
  citecolor=seplanblue,
  pdftitle={Guia da Secretaria — Orçamentos Temáticos},
  pdfauthor={SEPLAN — Governo do Acre},
}

\pagestyle{fancy}
\fancyhf{}
\fancyhead[L]{\small\textcolor{gray}{Guia da Secretaria — Orçamentos Temáticos}}
\fancyhead[R]{\small\textcolor{gray}{\thepage}}
\renewcommand{\headrulewidth}{0.4pt}

\titleformat{\section}{\Large\bfseries\color{seplanblue}}{\thesection}{1em}{}
\titleformat{\subsection}{\large\bfseries}{\thesubsection}{1em}{}
\titleformat{\subsubsection}{\normalsize\bfseries}{\thesubsubsection}{1em}{}

\setlist[itemize]{leftmargin=1.5em, topsep=4pt, itemsep=2pt}
\setlist[enumerate]{leftmargin=1.5em, topsep=4pt, itemsep=2pt}

\newtcolorbox{boxrepresentante}{
  breakable,
  colback=repgreen!8,
  colframe=repgreen!70,
  title={\textbf{Representante}},
  fonttitle=\bfseries\color{repgreen},
  left=8pt, right=8pt, top=6pt, bottom=6pt,
}

\newtcolorbox{boxrevisor}{
  breakable,
  colback=revblue!8,
  colframe=revblue!70,
  title={\textbf{Revisor interno}},
  fonttitle=\bfseries\color{revblue},
  left=8pt, right=8pt, top=6pt, bottom=6pt,
}

\newtcolorbox{boxatencao}{
  breakable,
  colback=warnamber!8,
  colframe=warnamber!70,
  title={\textbf{Atenção}},
  fonttitle=\bfseries\color{warnamber},
  left=8pt, right=8pt, top=6pt, bottom=6pt,
}

\newtcolorbox{boxdica}{
  breakable,
  colback=tipcyan!8,
  colframe=tipcyan!70,
  title={\textbf{Dica}},
  fonttitle=\bfseries\color{tipcyan},
  left=8pt, right=8pt, top=6pt, bottom=6pt,
}

\newcommand{\platform}[1]{\texttt{#1}}
\newcommand{\field}[1]{\textbf{#1}}

% Glossário — use tabela explícita no .tex (tabularx não funciona bem em \newenvironment)
